\scanpage{422}
\mspage{161}
une Hot-équivalence. D'autre part, $\unsure{\underline{\mathrm{Tz}}|\underline{\mathrm{Ty}}|\underline{\mathrm{Iz}}|\underline{\mathrm{Iy}}}$ ayant un objet final est asphérique, donc $X\times\underline{\mathrm{Tz}}\xrightarrow{\mathrm{pr}_1}X$ est une Hot-équivalence, d'où résulte aussitôt que $i$ est une Hot-équivalence. Donc de même $ji$, or $ji=\alpha_1$, qed.

2) Prouvons p. ex. la lissité de $\underline{\mathrm{Ch}}(X)\to X\times X$, la propreté s'en déduira en \deleted{\unsure{utilisant}}\inserted{appliquant} ce résultat \deleted{pour}\inserted{à} $X^{\mathrm{opp}}$. Soit donc
\[
\unsure{\backslash}(a',b')\longrightarrow(\unsure{a},\unsure{b})\qquad\text{i.e.}\quad a'\mathrel{\overset{u}{\unsure{\text{---}|\to}}}a,\ b'\xrightarrow{v}b
\]
une flèche dans $X\times X$, et $c$ un objet de $\underline{\mathrm{Ch}}(X)$ au dessus de $(a,b)$ :
\diagram{\begin{tikzcd}[ampersand replacement=\&,column sep=huge]
a'\arrow[r,dashed,"c'","{\rightarrow}"']\arrow[d,"u"'] \& b'\arrow[d,"v"]\\
a\arrow[r,dash,"{\substack{\rightarrow\\\deleted{\unsure{\leftarrow}}\\c}}"'] \& b.
\end{tikzcd}}
On regarde la catégorie $\mathcal C(c,u,v)$ des couples $(c',w)$ qui \unsure{relèvent|relivent|relient} $(c,(u,v))$, \deleted{avec \unsure{$\mathcal C$ cat.}}\inserted{où} \deleted{\unsure{de|dans}} $c'$ est un chemin (de type non précisé) de $a'$ vers $b'$, et \deleted{\unsure{avec}}
\[
w:c'\longrightarrow c
\]
une flèche dans $\underline{\mathrm{Ch}}(X)$ qui sur les \inserted{objets-}sources induise $u$, et sur les objets buts induise $v$ (ainsi $c'$ est un chemin de $a'$ vers $b'$). Il faut
% archive sheet number: 421
